\section{财富对选项价值的调制}\label{chap3_6:results_offer_value}
\ref{chap3_4:results_wealth_tuning}~节和\ref{chap3_5:results_wealth_updating}~节已确立前额叶皮层在单神经元和群体水平上对财富状态的稳定表征。在此基础上，我们进一步考察这一表征如何影响选项价值的编码进而参与经济决策的核心计算。根据边际效用递减原理，赚取价值应随财富增加而降低，消费价值应随财富增加而升高（图~\ref{fig5_1:example}A、C）；若前额叶神经元对价值的编码受财富调制，则其放电率应呈现类似的财富依赖模式。图~\ref{fig5_1:example}B和D展示的示例OFC神经元符合此预测：赚取选项引发的神经元反应随财富增加而下降，消费选项的反应则随财富增加而单调上升。

\begin{figure}[H]
    \centering
    \includegraphics[width=0.8\textwidth]{Figure5_1_Example}
    \bicaption{\enspace 编码选项价值的示例神经元}{\enspace Representative neurons encoding offer values}
    \label{fig5_1:example}
\end{figure}
\fignotetext{
    \textbf{(A) 效用模型对赚取价值的理论预测。}纵轴为模型估计的赚取选项主观价值，横轴为财富水平（五组），颜色深浅表示客观赚取速率（越深表示速率越高）。高赚取速率选项的价值整体高于低速率选项，且随着财富增加，赚取价值单调递减。\\
    \textbf{(B) 示例OFC神经元对赚取价值的编码。}纵轴为钱包更新期（$[-200, 300]$~毫秒）内的平均放电率（均值$\pm$标准误），其余布局同(A)。该神经元对赚取价值的编码方向与(A)中的模型预测一致。\\
    \textbf{(C) 效用模型对消费价值的理论预测。}纵轴为模型估计的消费选项主观价值，颜色深浅表示客观奖励量（越深表示奖励量越大）。随财富增加，消费价值单调递增，其余布局同(A)。\\
    \textbf{(D) 示例OFC神经元对消费价值的编码。}布局同(B)和(C)。该神经元对消费价值的编码方向与(C)中的模型预测一致。}
    {\textbf{(A) Model predictions for earning value.} The y-axis shows the model-estimated subjective earning value; the x-axis shows wealth level (five groups); color lightness encodes objective earning rate (darker = higher rate). Earn value decreases monotonically with wealth, and options with higher earning rates have higher value overall. \\
    \textbf{(B) Earning value encoding of a representative OFC neuron.} The y-axis shows mean firing rate (mean $\pm$ SEM) averaged over the wallet update epoch ($[-200, 300]$~ms); remaining layout as in (A). This neuron's firing rate is consistent with the pattern predicted in (A). \\
    \textbf{(C) Model predictions for spending value.} The y-axis shows model-estimated subjective spending value; color lightness encodes objective reward size (darker = larger reward). Spending value increases monotonically with wealth; remaining layout as in (A). \\
    \textbf{(D) Spending value encoding of a representative OFC neuron.} Layout as in (B). This neuron's firing rate is consistent with the pattern predicted in (C).}

由于自由$\alpha$模型导出的财富调制价值（$V_E$、$V_S$）与固定$\alpha$模型导出的财富无关价值（$V_E^{\mathrm{fixed}}$、$V_S^{\mathrm{fixed}}$）高度相关（Pearson相关系数：赚取价值在0.84-0.92之间，消费价值均超过0.99），直接在线性模型中同时纳入两者会产生共线性问题。在高度共线的情况下，回归系数估计变得不稳定——两个相关变量的系数可以任意分配，即使其中一个变量真实贡献为零，也可能被分配到非零系数。因此，我们无法仅凭回归系数判断单个神经元是否编码了被财富调制的价值。为绕过这一问题，我们转而采用联合编码分析（见~\ref{chap2:methods_neural_offer_joint}~节），检验单神经元对价值的编码方向与对财富的编码方向之间是否存在系统性的关联。该分析使用固定$\alpha$模型拟合的选项价值，而非自由$\alpha$模型，因为后者已包含财富调制信息，若作为回归变量，财富的调制效应可能被价值项吸收，导致联合编码中财富项的独立贡献被低估。联合编码分析的零假设是价值编码与财富编码在单神经元层面相互独立，同时对两者具有显著编码的神经元应随机分布于$2\times2$列联表四个格中。若财富确实调制了价值编码，则应观察到系统性偏离：消费价值编码与财富编码呈对角关联（富时消费价值更高），赚取价值编码与财富编码呈反对角关联（富时赚取价值更低）。Fisher精确检验显示，消费价值的对角模式在钱包更新期三个脑区均显著（图~\ref{fig5_2:joint}右列）；赚取价值与财富的反对角模式在三个脑区均未达显著（图~\ref{fig5_2:joint}左列），提示赚取价值的财富调制信号较弱，或需更敏感的方法检验。

\begin{figure}[H]
    \centering
    \includegraphics[width=\textwidth]{Figure5_2_joint}
    \bicaption{\enspace 选项价值与财富水平的联合编码}{\enspace Joint encoding of offer value and wealth}
    \label{fig5_2:joint}
\end{figure}
\fignotetext{
    \textbf{(A) 选项呈现期，赚取价值与财富的方向联合分布。}每格内部为$4\times3$排列：四行分别为聚合数据、猕猴W、猕猴U、猕猴T，三列分别为ACC、DLPFC、OFC。每个$2\times2$列联表展示同时编码选项价值和财富水平（$p < 0.05$）的神经元在各方向组合中的比例，采用左尾Fisher精确检验评估反对角集中程度（$V_E^{\mathrm{fixed}}+$/财富$-$ 与 $V_E^{\mathrm{fixed}}-$/财富$+$）。\\
    \textbf{(B) 选项呈现期，消费价值与财富的方向联合分布。}采用右尾Fisher精确检验评估对角集中程度（$V_S^{\mathrm{fixed}}+$/财富$+$ 与 $V_S^{\mathrm{fixed}}-$/财富$-$），其余布局同(A)。\\
    \textbf{(C) 钱包更新期，赚取价值与财富的方向联合分布。}布局同(A)。聚合数据中，$V_E^{\mathrm{fixed}}$与财富的反对角关联在三个脑区均未达显著（ACC：$p = 0.22$；DLPFC：$p = 0.25$；OFC：$p = 0.96$）。\\
    \textbf{(D) 钱包更新期，消费价值与财富的方向联合分布。}采用右尾Fisher精确检验评估正对角集中程度，其余布局同(C)。聚合数据中，$V_S^{\mathrm{fixed}}$与财富的对角关联在三个脑区均显著（ACC：$p = 9.4\times10^{-8}$；DLPFC：$p = 5.7\times10^{-7}$；OFC：$p = 4.0\times10^{-10}$）。}
    {\textbf{(A) Offer period, joint directional distribution of earning value and wealth.} Within each cell, panels are arranged in a $4\times3$ grid: four rows for pooled data, Monkey W, Monkey U, and Monkey T; three columns for ACC, DLPFC, and OFC. Each $2\times2$ contingency table shows the proportion of neurons with significant encoding of both offer value and wealth ($p < 0.05$) in each direction combination. A left-tailed Fisher exact test assesses anti-diagonal concentration ($V_E^{\mathrm{fixed}}+$/Wealth$-$ and $V_E^{\mathrm{fixed}}-$/Wealth$+$). \\
    \textbf{(B) Offer period, joint directional distribution of spending value and wealth.} A right-tailed Fisher exact test assesses diagonal concentration ($V_S^{\mathrm{fixed}}+$/Wealth$+$ and $V_S^{\mathrm{fixed}}-$/Wealth$-$); remaining layout as in (A). \\
    \textbf{(C) Token update epoch, joint directional distribution of earning value and wealth.} Layout as in (A). In the pooled data, the anti-diagonal pattern for earning value was non-significant in all three areas (ACC: $p = 0.22$; DLPFC: $p = 0.25$; OFC: $p = 0.96$). \\
    \textbf{(D) Token update epoch, joint directional distribution of spending value and wealth.} A right-tailed Fisher exact test assesses diagonal concentration; remaining layout as in (C). In the pooled data, the diagonal pattern for spending value was significant in all three areas (ACC: $p = 9.4\times10^{-8}$; DLPFC: $p = 5.7\times10^{-7}$; OFC: $p = 4.0\times10^{-10}$).}

我们进一步采用嵌套模型比较（参见~\ref{chap2:methods_neural_offer_time}~节），直接检验财富调制的赚取价值（$V_E$）能否在基线模型之上提供额外的解释力。基线模型（$M_1$）已包含不依赖财富的选项价值和财富水平，若在此基础上加入财富调制的赚取价值（$M_2$）能显著提升对神经元放电率的拟合能力，则构成财富调制存在的直接证据。此方法仅适用于赚取价值：赚取价值的财富水平是乘法交互（详见\ref{chap2:methods_BHV_utility}），在数学上无法被$M_1$中的独立项线性吸收，因而其增量贡献可被独立检验；而消费价值$V_S$与财富是加法交互，与$M_1$模型中已有回归量完全共线，无法通过增量方差分离，故后续分析仅针对赚取价值。结果显示，选项呈现期内的三个脑区$\Delta R^2$均未显著偏离0，但在钱包更新期内均显著为正（图~\ref{fig5_3:model}），提示财富对赚取价值的调制主要发生在结果反馈阶段。这一区别可能是由于选项呈现期的试次数更少导致的统计功效差异，为此，我们在钱包更新期内将有效试次下采样至与选项呈现期相同的数量，结果仍然显著，排除了试次数差异的影响。

联合编码分析以及嵌套模型比较的结果共同表明，财富水平系统地调制了前额叶皮层对选项价值的编码。

\begin{figure}[H]
    \centering
    \includegraphics[width=0.9\textwidth]{Figure5_3_ModelComparison}
    \bicaption{\enspace 财富调制效应的增量解释力}{\enspace Incremental variance explained by wealth modulation}
    \label{fig5_3:model}
\end{figure}
\fignotetext{四行分别为聚合数据、猕猴W、猕猴U、猕猴T；三列对应三个分析条件：选项呈现期（左）、试次数匹配的钱包更新期（中）、全部试次的钱包更新期（右）。中列与右列均分析钱包更新期，区别在于中列将有效试次下采样至与选项呈现期相同的数量，以排除试次数差异对统计显著性的影响。每张子图中，实线为$\Delta R^2 = R^2_{\mathrm{adj}}(M_2) - R^2_{\mathrm{adj}}(M_1)$在神经元间的均值$\pm$标准误，其中$M_1$为基线模型（包含$V_E^{\mathrm{fixed}}$、$V_S^{\mathrm{fixed}}$、$\tilde{W}_{24}$和选择），$M_2$在此基础上加入$V_E$；虚线为将$V_E$标签打乱后的零分布。采用调整后$R^2$以校正两模型自由参数数量的差异。横轴上方彩色三角标记$\Delta R^2$显著偏离零分布的时间窗（上三角为$\Delta R^2 > 0$，下三角为$\Delta R^2 < 0$；$p < 0.05$，双侧Wilcoxon符号秩检验；时间窗宽度为100毫秒）。聚合数据中，三个脑区在钱包更新期的$\Delta R^2$均显著为正，匹配试次数后依然成立；该效应在选项呈现期不显著。个体数据中，猕猴W在匹配试次的钱包更新期未达显著，但方向与聚合结果一致，使用全部试次后达到显著；其余两只猕猴的结果与聚合数据吻合。}{Four rows show data for the pooled sample, Monkey W, Monkey U, and Monkey T, respectively. Three columns correspond to three analysis conditions: offer onset period (left), trial-matched wallet update period (middle), and full-trial wallet update period (right). The middle and right columns both analyze the wallet update epoch; in the middle column, valid trials are subsampled to match the trial count of the offer onset period, controlling for the potential influence of trial count on statistical power. In each panel, the solid line shows the mean $\pm$ SEM of $\Delta R^2 = R^2_{\mathrm{adj}}(M_2) - R^2_{\mathrm{adj}}(M_1)$ across neurons, where $M_1$ is the baseline model (containing $V_E^{\mathrm{fixed}}$, $V_S^{\mathrm{fixed}}$,  $\tilde{W}_{24}$, and $Choice$) and $M_2$ additionally includes the $V_E$ term; the dashed line shows the null distribution obtained by permuting $V_E$ labels. Adjusted $R^2$ is used to correct for the difference in the number of free parameters between the two models. Colored triangles above the x-axis mark time windows in which $\Delta R^2$ significantly deviates from the null distribution (upward triangles: $\Delta R^2 > 0$; downward triangles: $\Delta R^2 < 0$; $p < 0.05$, two-sided Wilcoxon signed-rank test; time window width = 100 ms). In the pooled data, $\Delta R^2$ was significantly positive during the wallet update epoch in all three areas and remained significant after trial count matching; no significant effect was observed during the offer onset period. In individual monkeys, Monkey W did not reach significance in the trial-matched analysis but showed a consistent direction and reached significance with all trials; results for the other two monkeys were consistent with the pooled data.}